\section{Wildcard Matching}
\label{sec:dp-wildcard-matching}

\problemstatement{Given a string $s$ and a pattern $p$ containing
literal characters, \texttt{'?'} (matches any \textbf{single}
character), and \texttt{'*'} (matches any sequence of characters,
\textbf{including the empty sequence}), determine whether $p$ matches
$s$ \textbf{entirely}.}

\begin{patternbox}
This closing problem of the chapter is a deliberate capstone: it returns
to Section~\ref{sec:dp-lcs}'s match/mismatch two-pointer table, but the
``match'' condition is now governed by \textbf{pattern syntax}
(literal, \texttt{?}, or \texttt{*}) rather than simple character
equality, and \texttt{*} specifically introduces a branching choice
(match zero characters, or match one more) reminiscent of the unbounded
knapsack family's ``stay at the same index'' rule from earlier in this
book.
\end{patternbox}

\subsection*{Deriving the recurrence}

Let $f(i,j)$ be \texttt{true} if $s[0..i-1]$ (length $i$) matches
$p[0..j-1]$ (length $j$). Branch on $p[j-1]$:

\begin{itemize}
  \item If $p[j-1]$ is a literal character: it must equal $s[i-1]$
        exactly (so $i>0$ is required), and the rest must match:
        $f(i,j) = (i>0) \land (s[i{-}1]=p[j{-}1]) \land f(i{-}1,j{-}1)$.
  \item If $p[j-1] = \texttt{'?'}$: it matches any single character of
        $s$ (so $i>0$ is required, but no equality check):
        $f(i,j) = (i>0) \land f(i{-}1,j{-}1)$.
  \item If $p[j-1] = \texttt{'*'}$: it may match \textbf{zero}
        characters of $s$ (so the rest of $p$ must match all of $s$ so
        far: $f(i,j{-}1)$), \textbf{or} it may match \textbf{one more}
        character of $s$ while remaining available to match further
        characters too (so $i>0$ is required: $f(i{-}1,j)$). Either
        possibility is acceptable:
        $f(i,j) = f(i,j{-}1) \lor \big((i>0) \land f(i{-}1,j)\big)$.
\end{itemize}

Base cases: $f(0,0) = \texttt{true}$ (empty pattern matches empty
string). $f(0,j)$ for $j>0$: \texttt{true} only if $p[0..j-1]$ consists
\emph{entirely} of \texttt{'*'} characters (each can independently match
zero characters, collapsing the whole prefix to nothing).
$f(i,0)$ for $i>0$: always \texttt{false} (a non-empty string cannot be
matched by an empty pattern).

\subsection*{Approach 1: Recursion}

\begin{lstlisting}
#include <bits/stdc++.h>
using namespace std;

bool wildcardRecursive(int i, int j, string& s, string& p) {
    if (i == 0 && j == 0) return true;
    if (j == 0) return false; // i > 0 here
    if (i == 0) {
        // s is empty: only matches if remaining pattern is all '*'
        for (int k = 0; k < j; k++) if (p[k] != '*') return false;
        return true;
    }

    if (p[j - 1] == '*') {
        return wildcardRecursive(i, j - 1, s, p) ||
               wildcardRecursive(i - 1, j, s, p);
    }
    if (p[j - 1] == '?' || p[j - 1] == s[i - 1]) {
        return wildcardRecursive(i - 1, j - 1, s, p);
    }
    return false;
}

int main() {
    string s = "aa", p = "*";
    cout << (wildcardRecursive(s.size(), p.size(), s, p) ? "true" : "false") << endl;
    return 0;
}
\end{lstlisting}

\begin{complexitybox}[Approach 1 Complexity]
\textbf{Time.} $O(2^{n+m})$ in the worst case -- two branches per
\texttt{'*'} encountered.
\textbf{Space.} $O(n+m)$ recursion stack.
\end{complexitybox}

\subsection*{Approach 2: Memoization}

\begin{lstlisting}
#include <bits/stdc++.h>
using namespace std;

int wildcardMemo(int i, int j, string& s, string& p,
                  vector<vector<int>>& memo) {
    if (i == 0 && j == 0) return 1;
    if (j == 0) return 0;
    if (i == 0) {
        for (int k = 0; k < j; k++) if (p[k] != '*') return 0;
        return 1;
    }
    if (memo[i][j] != -1) return memo[i][j];

    bool result;
    if (p[j - 1] == '*') {
        result = wildcardMemo(i, j - 1, s, p, memo) ||
                  wildcardMemo(i - 1, j, s, p, memo);
    } else if (p[j - 1] == '?' || p[j - 1] == s[i - 1]) {
        result = wildcardMemo(i - 1, j - 1, s, p, memo);
    } else {
        result = false;
    }
    return memo[i][j] = result;
}

int main() {
    string s = "adceb", p = "*a*b";
    int n = s.size(), m = p.size();
    vector<vector<int>> memo(n + 1, vector<int>(m + 1, -1));
    cout << (wildcardMemo(n, m, s, p, memo) ? "true" : "false") << endl; // true
    return 0;
}
\end{lstlisting}

\begin{complexitybox}[Approach 2 Complexity]
\textbf{Time.} $O(n \cdot m)$.
\textbf{Space.} $O(n \cdot m)$ memo $+$ $O(n+m)$ recursion stack.
\end{complexitybox}

\subsection*{Approach 3: Tabulation}

\begin{lstlisting}
#include <bits/stdc++.h>
using namespace std;

bool wildcardTabulation(string& s, string& p) {
    int n = s.size(), m = p.size();
    vector<vector<bool>> dp(n + 1, vector<bool>(m + 1, false));
    dp[0][0] = true;

    for (int j = 1; j <= m; j++) {
        bool allStars = true;
        for (int k = 0; k < j; k++) if (p[k] != '*') { allStars = false; break; }
        dp[0][j] = allStars;
    }

    for (int i = 1; i <= n; i++) {
        for (int j = 1; j <= m; j++) {
            if (p[j - 1] == '*') {
                dp[i][j] = dp[i][j - 1] || dp[i - 1][j];
            } else if (p[j - 1] == '?' || p[j - 1] == s[i - 1]) {
                dp[i][j] = dp[i - 1][j - 1];
            } else {
                dp[i][j] = false;
            }
        }
    }
    return dp[n][m];
}

int main() {
    string s = "adceb", p = "*a*b";
    cout << (wildcardTabulation(s, p) ? "true" : "false") << endl; // true
    return 0;
}
\end{lstlisting}

\subsection*{Dry run of the tabulation table}

\dryrun{Take $s = \texttt{"adceb"}$, $p = \texttt{"*a*b"}$.}

Row $0$ ($s$ empty): $\textit{dp}[0][0]=\text{T}$;
$\textit{dp}[0][1]=\text{T}$ (\texttt{"*"} matches empty);
$\textit{dp}[0][2..4]=\text{F}$ (prefix \texttt{"*a"} etc.\ contain a
non-\texttt{*} character, so cannot match the empty string). Filling
forward: the first \texttt{*} absorbs \texttt{"adc"}, matching the
literal \texttt{a} against the \texttt{a} in \texttt{"adceb"}, the
second \texttt{*} absorbs nothing further, and the literal \texttt{b}
matches the trailing \texttt{b}. The final cell
$\textit{dp}[5][4]=\text{T}$, confirming the match.

\begin{complexitybox}[Approach 3 Complexity]
\textbf{Time.} $O(n \cdot m)$.
\textbf{Space.} $O(n \cdot m)$ table, no recursion stack.
\end{complexitybox}

\subsection*{Approach 4: Space Optimization}

\begin{lstlisting}
#include <bits/stdc++.h>
using namespace std;

bool wildcardSpaceOptimized(string& s, string& p) {
    int n = s.size(), m = p.size();
    vector<bool> prev(m + 1, false), current(m + 1, false);
    prev[0] = true;

    for (int j = 1; j <= m; j++) {
        bool allStars = true;
        for (int k = 0; k < j; k++) if (p[k] != '*') { allStars = false; break; }
        prev[j] = allStars;
    }

    for (int i = 1; i <= n; i++) {
        current[0] = false;
        for (int j = 1; j <= m; j++) {
            if (p[j - 1] == '*') {
                current[j] = current[j - 1] || prev[j];
            } else if (p[j - 1] == '?' || p[j - 1] == s[i - 1]) {
                current[j] = prev[j - 1];
            } else {
                current[j] = false;
            }
        }
        prev = current;
    }
    return prev[m];
}

int main() {
    string s = "adceb", p = "*a*b";
    cout << (wildcardSpaceOptimized(s, p) ? "true" : "false") << endl; // true
    return 0;
}
\end{lstlisting}

\begin{complexitybox}[Approach 4 Complexity]
\textbf{Time.} $O(n \cdot m)$.
\textbf{Space.} $O(m)$ for two rolling rows.
\end{complexitybox}

\begin{takeawaybox}
This chapter closes by combining two of its biggest ideas: the
match/mismatch two-pointer table from Section~\ref{sec:dp-lcs}, and a
branching ``match zero or match one more'' choice for \texttt{*}
strongly reminiscent of the unbounded-knapsack reuse rule from the DP on
Subsequences chapter. Recognizing that a special pattern character
introduces exactly this same kind of ``stay and try again'' branch is
the key insight that unlocks Wildcard Matching (and its sibling,
Regular Expression Matching, solved by an almost identical table) without
needing an entirely new derivation strategy.
\end{takeawaybox}
